Per un fil recte molt llarg circula un corrent d'1,5~A en el sentit positiu de la direcció $y$, seguint la línia $x = -3,0\ \mathrm{cm}$. Un altre fil amb les mateixes característiques, pel qual també circula un corrent d'1,5~A en el sentit positiu de la direcció $y$, segueix la línia $x = 3,0\ \mathrm{cm}$, com mostra la figura.

\figura[0.35]{fig1.png}

\begin{apartat}{1}
Calculeu el camp magnètic (mòdul, direcció i sentit) en $x = 0$ i feu un esquema que justifiqui el resultat.
\end{apartat}

\begin{apartat}{1}
Calculeu el camp magnètic (mòdul, direcció i sentit) en $x = 5,0\ \mathrm{cm}$ i feu un esquema que justifiqui el resultat.
\end{apartat}

\blocetiquetat{Dada:}{$\mu_0 = 4\pi\times 10^{-7}\ \mathrm{T\,m\,A^{-1}}$.}
\nota{El mòdul del camp magnètic creat per un fil conductor infinit pel qual circula una intensitat de corrent $I$ és: $B = \dfrac{\mu_0 I}{2\pi r}$, en què $r$ és la distància al fil conductor.}
